% Statement and cited-input packet. Target proof is separately embargoed.
% Frozen II source SHA256: 74497bfe327fae25c9cdcb17334b3bf42a9793a9df7c0a4404d97568f2f760f0
% Target owner: cdp:st:C4c; tag: 01FM

\begin{proposition}[The finite singular-fibre test and a fixed unavoidable fibre]\hypertarget{tag-01FM}{}\label{cdp:st:C4c}


Let $f:Y\to\PP^1$ be a proper connected-fibre map from a compact complex threefold with $b_1(Y)=b_2(Y)=0$ and smooth general fibre a complex two-torus. Let $\Sigma$ be its finite set of critical values, and fix a connected component $P$ of $\Pic(Y)$. For $c\in\Sigma$ retain the entire Cartier fibre $F_c$, including its multiplicities, and put
\[
 Z_c=\{L\in P:H^0(F_c,\Omega_Y^1|_{F_c}\otimes
                  (K_Y\otimes L^{-1})|_{F_c})\ne0\}.
                                                               \tag{C.18}\hypertarget{tag-008W}{}\label{cdp:eq:C.18}
\]
Then the following are equivalent: every $Z_c$ is a proper subset of $P$; and $H^2(Y,T_Y\otimes L)=0$ on a dense open subset of $P$. In that case the generic twist in $P$ satisfies
\[
 H^q(Y,T_Y\otimes L)=0\quad(q\ne1),\qquad
 \chi(Y,T_Y\otimes L)\le0,\qquad c_3(Y)\le0.          \tag{C.19}\hypertarget{tag-008X}{}\label{cdp:eq:C.19}
\]
For a reduced fibre, membership in the complement of $Z_c$ is exactly the two explicitly typed kernel-and-connecting-map conditions of \ref{cdp:st:A.1.1}, with kernel $N^*\otimes\mathcal J^{-1}$ computed in A.2. For $F_c=mS$, the equivalent complete test is the finite lifting filtration of \ref{cdp:st:A.3.1}. These tests use only sheaves and maps on the finitely many singular full fibres.

When $b_3(Y)=0$, at least one fixed critical fibre $F_c$ satisfies $Z_c=P$. Its support is irreducible, by \ref{cdp:st:A.3.5}. Under the integral assumption $H^2(Y,\Z)=0$, the Picard group is connected, so this is a single fixed fibre for which the ambient group in \eqref{cdp:eq:C.18} is nonzero for every line bundle on $Y$.
\end{proposition}

% Cited input cdp:st:A.1.1; statement or labelled formula only.
\begin{proposition}[The conormal filtration and its exact section criterion]\hypertarget{tag-01AD}{}\label{cdp:st:A.1.1}

\label{wbII:prop:criterion}\label{wbII:prop:filtration}\label{wbII:lem:conormal}\label{wbII:lem:conormal-injective}\label{wbII:prop:missing-global-section-criterion}\label{wbII:prop:missing-cotangent-filtration}
The conormal map $d:N^*\to V$, defined by $[F]\mapsto dF|_S$ for a local defining equation $F$, is injective. With $K=\ker(V\to Q)$, there are canonical exact sequences
\begin{align}
0&\longrightarrow N^*\xrightarrow{d}K\longrightarrow T\longrightarrow0,
\hypertarget{tag-003W}{}\label{cdp:A:filtration1}\\
0&\longrightarrow K\longrightarrow V\longrightarrow Q\longrightarrow0.
\hypertarget{tag-003X}{}\label{cdp:A:filtration2}
\end{align}
Let $\partial_A:H^0(S,Q\otimes A)\to H^1(S,K\otimes A)$ and $\delta_A:H^0(S,T\otimes A)\to H^1(S,N^*\otimes A)$ be their connecting maps. Then
\[
 H^0(S,V\otimes A)=0
 \quad\Longleftrightarrow\quad
 H^0(S,K\otimes A)=0\ \hbox{and}\ \partial_A\hbox{ is injective}.
\tag{A.1}\hypertarget{tag-003Y}{}\label{cdp:eq:A.1}
\]
Under the explicit hypotheses $H^0(S,N^*\otimes A)=0$ and $H^0(S,Q\otimes A)=0$, $H^0(S,V\otimes A)$ is canonically $\ker\delta_A$.
\end{proposition}

% Cited input cdp:st:A.3.1; statement or labelled formula only.
\begin{proposition}[All thickenings and their lifting maps]\hypertarget{tag-01AL}{}\label{cdp:st:A.3.1}


Let $S$ be any effective Cartier divisor in a complex manifold, let $E$ be any locally free sheaf on that manifold, and let $E_k=E|_{kS}$. For each $k\ge2$ there is a canonical exact sequence
\[
0\longrightarrow (E|_S)\otimes(N^*)^{k-1}
 \longrightarrow E_k\longrightarrow E_{k-1}\longrightarrow0.
\tag{A.10}\hypertarget{tag-003Z}{}\label{cdp:eq:A.10}
\]
Write $G_{k-1}=(E|_S)\otimes(N^*)^{k-1}$ and let
$\epsilon_k:H^0((k-1)S,E_{k-1})\to H^1(S,G_{k-1})$
be the connecting map. Then
\[
0\longrightarrow H^0(S,G_{k-1})\longrightarrow H^0(kS,E_k)
 \longrightarrow\ker\epsilon_k\longrightarrow0.
\tag{A.11}\hypertarget{tag-0040}{}\label{cdp:eq:A.11}
\]
In particular $H^0(kS,E_k)=0$ exactly when $H^0(S,G_{k-1})=0$ and $\epsilon_k$ is injective. Starting with $H^0(S,E|_S)$, equations \eqref{cdp:eq:A.10}--\eqref{cdp:eq:A.11} compute the group on $mS$ through all its extension maps. The vanishing of every $H^0(S,G_j)$, $0\le j<m$, is a sufficient criterion. These statements apply to $E=\Omega_Y^1\otimes B$ for every line bundle $B$, and include smooth, normal-crossings, and all other reduced supports.
\end{proposition}

% Cited input cdp:st:A.3.5; statement or labelled formula only.
\begin{proposition}[Every reducible full fibre has generic ambient vanishing]\hypertarget{tag-01AP}{}\label{cdp:st:A.3.5}


Let $Y$ be a smooth compact complex threefold with $H^2(Y,\Q)=0$, and let $F=\sum_i m_iS_i$ be an effective Cartier divisor with connected reducible support. Every connected component of $\Pic(Y)$ contains a dense open subset, with proper closed analytic complement, on which
\[
 H^0(F,(\Omega_Y^1\otimes L)|_F)=0.
 \tag{A.15b}\hypertarget{tag-0046}{}\label{cdp:eq:A.15b}
\]
The multiplicities may be unequal. For a proper connected fibration over a compact curve, the vanishing is simultaneous on all reducible full fibres after intersecting finitely many such open sets. Consequently a full fibre with ambient sections for every Picard twist has irreducible support.
\end{proposition}

% Cited input cdp:st:D.1.11; statement or labelled formula only.
\begin{proposition}\hypertarget{tag-01FW}{}\label{cdp:st:D.1.11}


Let $f:Y\to\PP^1$ be a proper connected-fibre holomorphic map from
a connected compact complex manifold with $b_1(Y)=b_2(Y)=0$.
In every component of $\Pic(Y)$, the line
bundles which are non-torsion on every full fibre form a dense
$G_\delta$ set of full Lebesgue measure in its complex vector-space
coordinates. If evaluation of global sections of $R^1f_*\OO_Y$
has rank one, then $\Pic^0(Y)=\C$ and the exceptional set is
countable. For two-torus fibres this alternative holds whenever
the splitting has a negative summand.
In particular the general-line-bundle conclusion of CDP20
Proposition 5.6 holds with this precise quantifier, without the
nonnegative-summand assertion of Proposition 5.5.
\end{proposition}

% Cited input cdp:eq:A.10; statement or labelled formula only.
\begin{proposition}[All thickenings and their lifting maps]\hypertarget{tag-01AL}{}\label{cdp:st:A.3.1}


Let $S$ be any effective Cartier divisor in a complex manifold, let $E$ be any locally free sheaf on that manifold, and let $E_k=E|_{kS}$. For each $k\ge2$ there is a canonical exact sequence
\[
0\longrightarrow (E|_S)\otimes(N^*)^{k-1}
 \longrightarrow E_k\longrightarrow E_{k-1}\longrightarrow0.
\tag{A.10}\hypertarget{tag-003Z}{}\label{cdp:eq:A.10}
\]
Write $G_{k-1}=(E|_S)\otimes(N^*)^{k-1}$ and let
$\epsilon_k:H^0((k-1)S,E_{k-1})\to H^1(S,G_{k-1})$
be the connecting map. Then
\[
0\longrightarrow H^0(S,G_{k-1})\longrightarrow H^0(kS,E_k)
 \longrightarrow\ker\epsilon_k\longrightarrow0.
\tag{A.11}\hypertarget{tag-0040}{}\label{cdp:eq:A.11}
\]
In particular $H^0(kS,E_k)=0$ exactly when $H^0(S,G_{k-1})=0$ and $\epsilon_k$ is injective. Starting with $H^0(S,E|_S)$, equations \eqref{cdp:eq:A.10}--\eqref{cdp:eq:A.11} compute the group on $mS$ through all its extension maps. The vanishing of every $H^0(S,G_j)$, $0\le j<m$, is a sufficient criterion. These statements apply to $E=\Omega_Y^1\otimes B$ for every line bundle $B$, and include smooth, normal-crossings, and all other reduced supports.
\end{proposition}

% Cited input cdp:eq:A.11; statement or labelled formula only.
\begin{proposition}[All thickenings and their lifting maps]\hypertarget{tag-01AL}{}\label{cdp:st:A.3.1}


Let $S$ be any effective Cartier divisor in a complex manifold, let $E$ be any locally free sheaf on that manifold, and let $E_k=E|_{kS}$. For each $k\ge2$ there is a canonical exact sequence
\[
0\longrightarrow (E|_S)\otimes(N^*)^{k-1}
 \longrightarrow E_k\longrightarrow E_{k-1}\longrightarrow0.
\tag{A.10}\hypertarget{tag-003Z}{}\label{cdp:eq:A.10}
\]
Write $G_{k-1}=(E|_S)\otimes(N^*)^{k-1}$ and let
$\epsilon_k:H^0((k-1)S,E_{k-1})\to H^1(S,G_{k-1})$
be the connecting map. Then
\[
0\longrightarrow H^0(S,G_{k-1})\longrightarrow H^0(kS,E_k)
 \longrightarrow\ker\epsilon_k\longrightarrow0.
\tag{A.11}\hypertarget{tag-0040}{}\label{cdp:eq:A.11}
\]
In particular $H^0(kS,E_k)=0$ exactly when $H^0(S,G_{k-1})=0$ and $\epsilon_k$ is injective. Starting with $H^0(S,E|_S)$, equations \eqref{cdp:eq:A.10}--\eqref{cdp:eq:A.11} compute the group on $mS$ through all its extension maps. The vanishing of every $H^0(S,G_j)$, $0\le j<m$, is a sufficient criterion. These statements apply to $E=\Omega_Y^1\otimes B$ for every line bundle $B$, and include smooth, normal-crossings, and all other reduced supports.
\end{proposition}

